\chapter*{Future Work}
\chaptermark{Future Work}
\addcontentsline{toc}{chapter}{Future Work}  
\label{ch:future-work}

While the proposed framework produces strong quantitative results and qualitatively 
sensible behaviour, several structural limitations constrain the current design. These 
are discussed per component below, together with the most natural directions for 
future work.

\section{User-Side Recommender}

The \gls{TT} model operates on implicit feedback: application events are treated as 
binary positive signals, meaning \emph{applied to} is assumed to be a \emph{good fit} 
and \emph{not applied to} a \emph{poor fit}. A user may apply speculatively to 
positions for which they are underqualified, or conversely fail to apply to a highly 
relevant job due to limited visibility, introducing label noise. Two complementary 
improvements to the training objective address this. First, the loss function could be 
adapted to correct for false negatives: ~\cite{yin2026correctweightsimpleeffective} 
proposes a Corrected and Weighted (CW) loss that re-calibrates the assumed negative 
distribution using Positive-Unlabeled learning and dynamically down-weights uncertain 
negatives that are likely mislabelled. Second, even with corrected labels, Binary 
Cross-Entropy treats each example independently while the actual task is a ranking 
problem. Replacing it with a ranking-aware objective such as Bayesian Personalised 
Ranking~\cite{rendle2012bprbayesianpersonalizedranking} would align training more 
directly with the evaluation metrics without any architectural change.

Finally, the model optimises purely for user preference with no explicit notion of 
fairness or congestion, leaving the market-maker framework as a post-processing 
correction. A more principled direction would be to incorporate the congestion 
objective directly into training, so that the model learns fairer representations from 
the start. The work of~\cite{mashayekhi2023reconreducingcongestionjob} demonstrates 
this by combining an optimal transport component with a job recommendation model in a 
joint optimisation problem, showing that congestion can be reduced without sacrificing 
recommendation quality.

\section{Recruiter-Side Recommender}

The biggest limitation is that it learns to replicate the judgments of an \gls{LLM} rather than actual human recruiters, using 
features that only partially capture what the LLM itself reasons over. On the label 
side, \glspl{LLM} may carry biases inherited from pretraining 
data~\cite{iso2025evaluatingbiasllmsjobresume}, and access to real recruiter feedback 
would remove this layer of indirection entirely. On the feature side, the current 
heuristic and semantic signals do not fully capture the nuances of job descriptions 
and candidate profiles. More expressive feature engineering or even a more complex model architecture could improve the model's ability to predict recruiter judgments, though this would require more data. 

A separate limitation is the binary classification approach. The four-level \gls{LLM} 
fit judgments (\textit{Very Good Fit}, \textit{Good Fit}, \textit{Poor Fit}, 
\textit{Very Poor Fit}) are mapped to binary labels before training, removing a 
meaningful distinction. Training directly on the ordinal fit levels using ordinal 
regression~\cite{wang2025surveyordinalregressionapplications} would preserve this 
ordinal structure and give the model a better signal, potentially producing 
better-calibrated suitability scores at inference time.

\section{Congestion Balancing}

The first limitation is the manual setting of all threshold and 
shape parameters ($\tau_{\min}$, $\tau_{\max}$, $H$, $B$, $P$, and the market maker 
scale parameters), which are pure design choices.
Second, the penalty functions themselves are static, applying the same shape 
regardless of whether a job is a popular or not. Framing 
parameter selection as an optimisation problem, and allowing penalty shapes to 
adapt to entity-level characteristics, would address both limitations simultaneously.

The final limitation is structural which is the order dependence of the sequential simulation. Users 
and job inboxes are processed one by one, and the running counters $c_j^{(t)}$ and 
$v_u^{(t)}$ accumulate as the simulation progresses, meaning early-processed users 
shape the penalty landscape seen by later ones. This means for example that a user at the beginning gets better recommended jobs than a user towards the end, which is unfair. Future work could examine the impact of processing order an solution this by for example randomising the order at each iteration or processing users in parallel.
